\subsection{Token Consumption Analysis}
\label{subsec:token-usage}

While repair success measures functional correctness, it overlooks a critical operational dimension: token consumption. We analyze token consumption patterns across agents to understand the computational cost of repair attempts and whether successful repairs differ from failed ones in token consumption.

We measure API token consumption in two dimensions. Input tokens represent the total tokens sent to the model across all repair attempts for a bug. Output tokens represent the total tokens generated by the model across all repair attempts for a bug. 
% We exclude bugs with processing duration less than 30 seconds, as these represent incomplete or aborted attempts. Token counts are aggregated per bug and stratified by repair outcome (successful vs failed).

\header{Token Consumption Across Agents}

Table~\ref{tab:token-usage-statistics} presents token consumption statistics across all four agents. \claudecode demonstrates the lowest per-bug input token consumption, representing 63$\times$ less than \codex (61.3K), 1,372$\times$ less than \geminicli (1.33M), and 1,753$\times$ less than \qwencode (1.70M). Prompt caching enables this reduction by reusing previously cached project context. 
% Other agents do not use caching and must resend all context (bug description, test cases, source code, dependencies) with every API call.

For output tokens per bug, \codex generates the fewest at 1.6K tokens on average, followed by \geminicli at 4.2K, \qwencode at 7.1K, and \claudecode at 18.1K.

Figure~\ref{fig:token-distribution-input} and Figure~\ref{fig:token-distribution-output} show violin plots of token distributions by outcome. Figure~\ref{fig:token-distribution-input} shows new input tokens only for \claudecode (not total input including cache).

\claudecode's prompt caching divides input tokens into three distinct categories. New input tokens represent uncached prompt content unique to each request, such as bug-specific information, test failures, and iterative repair instructions (successful repairs average 969 per bug). Cache creation tokens represent prompt content written to cache for reuse across multiple API calls, such as project source code, dependency files, and static context (successful repairs average 147K per bug). Cache read tokens represent previously cached content retrieved in subsequent API invocation (successful repairs average 3.70M per bug). Figure~\ref{fig:claude-input-token-distribution-violin} reveals that cache read tokens dominate at 96\% of total input volume. 
% Cache reads cost 90\% less than regular input (\$0.30 vs \$3.00 per million), making repeated context retrieval economical despite high token volumes.


\header{Token Consumption by Repair Outcome}

Failed repairs exhibit distinct token consumption patterns across agents, as shown by wider failure distributions (red violins) in Figure~\ref{fig:token-distribution-input} and Figure~\ref{fig:token-distribution-output}. \qwencode's failed repairs use 39.3\% more input tokens and 35.0\% more output tokens compared to successful repairs. \geminicli exhibits the most dramatic increase, with failed repairs consuming 343.2\% more input tokens and 190.5\% more output tokens. \codex shows 79.4\% more input tokens and 75.0\% more output tokens for failures. These patterns indicate that failures do not terminate early; agents persist through multiple unsuccessful attempts before giving up.

\claudecode demonstrates a different pattern due to prompt caching. Failed repairs consume 68.2\% fewer new input tokens (308 vs 969) but 77.1\% more output tokens (32.0K vs 18.1K). The reduced new input suggests failures trigger fewer repair iterations. However, Table~\ref{tab:claude-token-usage} shows failures consume 143.7\% more cache creation tokens and 152.4\% more cache read tokens than successful repairs, indicating extensive context retrieval despite fewer new requests.

\header{Implications for Deployment}

Token consumption directly impacts operational cost in production deployments. \sonnet pricing includes \$3 per million input tokens, \$3.75 per million cache creation tokens, \$0.30 per million cache read tokens, and \$15 per million output tokens. For 288 bugs processed by \claudecode, total API costs are approximately \$598 (\$2.08 per bug). For 372 bugs processed by \codex, costs are approximately \$94 (\$0.25 per bug). \claudecode attains the highest repair rate (94.8\%) while \codex demonstrates 87.1\% repair rate, resulting in an 8.3$\times$ cost difference per bug.

The resource premium for failed repairs (up to 343\% more tokens) represents a hidden tax on agents with low repair rates. \qwencode's 28.7\% repair rate means 71.3\% of bugs consume excess resources without producing valid fixes. \geminicli's 41.3\% repair rate similarly incurs high failure costs. This failure tax amplifies operational costs for low-performing agents, making them economically unsuitable.

Organizations must balance repair accuracy against operational cost. For cost-sensitive deployments, \codex provides strong repair rates at minimal cost. For quality-critical applications where maximizing repair success justifies higher expenses, \claudecode offers superior accuracy. \geminicli and \qwencode impose high costs with low success rates, making them unsuitable for production deployment.








\begin{table}[t]
\centering
\caption{Token Consumption per bug on multi-hunk bugs}
\label{tab:token-usage-statistics}
\begin{tabular}{@{}lrrrr@{}}
\toprule
& \multicolumn{2}{c}{\textbf{Pass}} & \multicolumn{2}{c}{\textbf{Fail}} \\
\cmidrule(lr){2-3} \cmidrule(lr){4-5}
\textbf{Agent} & \textbf{Input} & \textbf{Output} & \textbf{Input} & \textbf{Output} \\
\midrule
\qwencode & 1.70M & 7.1K & 2.37M & 9.5K \\
\geminicli & 1.33M & 4.2K & 5.90M & 12.1K \\
\codex & 61.3K & 1.6K & 110.0K & 2.7K \\
\claudecode & 969 & 18.1K & 308 & 32.0K \\
\bottomrule
\end{tabular}
\begin{tablenotes}
\small
\item All values are means per bug.
\item \claudecode shows new input tokens only; cache tokens detailed in Table~\ref{tab:claude-token-usage}.
\item Bugs with duration $<$ 30 seconds excluded (incomplete processing).
\end{tablenotes}
\end{table}

\begin{table}[t]
\centering
\caption{Claude Code Token Usage with Prompt Caching Breakdown}
\label{tab:claude-token-usage}
\begin{tabular}{@{}lrrrrr@{}}
\toprule
\textbf{Outcome} & \textbf{n} & \textbf{New Input} & \textbf{Cache Creation} & \textbf{Cache Read} & \textbf{Total Input} \\
\midrule
Pass & 273 & 969 & 147K & 3.70M & 3.85M \\
Fail & 15  & 308 & 359K & 9.34M & 9.70M \\
\bottomrule
\end{tabular}
\begin{tablenotes}
\small
\item All values are means per bug. Total Input = New Input + Cache Creation + Cache Read.
\item New Input: uncached tokens sent with each request (\$3.00 per M).
\item Cache Creation: tokens written to cache (\$3.75 per M, 25\% premium).
\item Cache Read: tokens retrieved from cache (\$0.30 per M, 90\% discount).
\item Failed repairs consume 152\% more total input tokens than successful repairs.
\end{tablenotes}
\end{table}


\begin{figure}[t]
\centering
\includegraphics[width=\columnwidth]{analysis/figures/agent_input_token_distribution_violin.pdf}
\caption{Violin plot distributions of per-bug API input token consumption across four coding agents. Each data point represents one bug. Green violins represent successful repairs, red violins represent failed repairs. Each agent has independent y-axis scaling. \claudecode uses logarithmic scale due to extreme outliers. For \claudecode, this shows only new input tokens (not cache tokens); see Figure~\ref{fig:claude-input-token-distribution-violin} for complete breakdown.}
\label{fig:token-distribution-input}
\end{figure}

\begin{figure}[t]
\centering
\includegraphics[width=\columnwidth]{analysis/figures/agent_output_token_distribution_violin.pdf}
\caption{Violin plot distributions of per-bug API output token generation across four coding agents. Each data point represents one bug. Green violins represent successful repairs, red violins represent failed repairs. Each agent has independent y-axis scaling. Failed repairs generate more output tokens across all agents, indicating prolonged unsuccessful exploration before termination.}
\label{fig:token-distribution-output}
\end{figure}

\begin{figure}[t]
\centering
\includegraphics[width=\columnwidth]{analysis/figures/claude_input_token_distribution_violin.pdf}
\caption{\claudecode per-bug input token distribution by type and repair outcome. Each data point represents one bug. Three subplots show distributions for new input tokens, cache creation tokens, and cache read tokens. Green violins represent successful repairs, red violins represent failed repairs. Cache read tokens dominate at 96\% of total input volume. Failed repairs consume 152\% more total input tokens than successful repairs (9.70M vs 3.85M).}
\label{fig:claude-input-token-distribution-violin}
\end{figure}
